\documentclass[11pt,a4paper]{ctexart}

% === 页面设置 ===
\usepackage[top=1.5cm, bottom=1.5cm, left=1.8cm, right=1.8cm]{geometry}
\usepackage{enumitem}
\setlist{nosep, leftmargin=*}

% === 数学 ===
\usepackage{amsmath, amssymb, amsthm, bm}
\usepackage{mathrsfs}

% === 其他 ===
\usepackage{xcolor}
\usepackage{titlesec}
\usepackage{fancyhdr}
\usepackage{hyperref}

\hypersetup{colorlinks=true, linkcolor=blue, urlcolor=blue}

% === 标题格式 ===
\titleformat{\section}{\Large\bfseries}{}{0em}{}
\titlespacing{\section}{0pt}{8pt}{4pt}
\titleformat{\subsection}{\large\bfseries}{}{0em}{}
\titlespacing{\subsection}{0pt}{6pt}{2pt}

\pagestyle{fancy}
\fancyhf{}
\fancyhead[L]{\small 高级算法 · 考场速写版}
\fancyhead[R]{\small 18题 · 只抽2题}
\fancyfoot[C]{\thepage}
\renewcommand{\headrulewidth}{0.4pt}

% === 自定义命令 ===
\newcommand{\R}{\mathbb{R}}
\newcommand{\E}{\mathbb{E}}
\newcommand{\Var}{\operatorname{Var}}
\newcommand{\ind}[1]{\mathbf{1}_{[#1]}}
\newcommand{\norm}[1]{\|#1\|}
\newcommand{\abs}[1]{|#1|}
\newcommand{\inner}[1]{\langle #1 \rangle}
\newcommand{\calN}{\mathcal{N}}

\newenvironment{answer}{\begin{quote}\small}{\end{quote}}

\begin{document}

\title{\textbf{高级算法期末考试题库 \\ 考场速写版}}
\author{18题全 · Ch4降维9题 + Ch5流与LP 9题 · 只抽2题}
\date{\today}
\maketitle
\thispagestyle{fancy}

\setcounter{page}{1}

% ============================================================
\section*{Ch4：降维（Dimension Reduction）}
% ============================================================

% --- Q1 ---
\subsection*{Q1 \quad 低 $\ell_\infty$-norm 对高斯 JL 变换没有帮助}

\textbf{1.} $\Pi \in \R^{m\times d}$，$\pi_{ij} \sim \calN(0, 1/m)$ 独立。$\|x\|_2 = 1$。

$(\Pi x)_i = \sum_{j=1}^d \pi_{ij} x_j$。由正态分布可加性：
\[
\sum_j x_j \cdot \calN(0, \tfrac{1}{m}) \sim \calN\!\left(0,\; \tfrac{1}{m}\sum_j x_j^2\right) = \calN(0, \tfrac{1}{m})
\]
各坐标独立（不同 $i$ 的 $\pi_{ij}$ 无关）。\hfill $\square$

\textbf{2.} $\|\Pi x\|_2^2 = \sum_{i=1}^m (\Pi x)_i^2$。由(1)知 $(\Pi x)_i \sim \calN(0, 1/m)$ i.i.d.

令 $Z_i = \sqrt{m}\,(\Pi x)_i \sim \calN(0,1)$ i.i.d.，则
\[
\|\Pi x\|_2^2 = \frac{1}{m}\sum_{i=1}^m Z_i^2 \sim \frac{1}{m}\,\chi^2_m
\]
该分布仅依赖于 $m$，与 $x$ 的坐标值（包括 $\|x\|_\infty$）无关。\hfill $\square$

\textbf{3.} 高斯 JL 的集中不等式：
\[
\Pr\!\big[\,\big|\|\Pi x\|_2^2 - 1\big| > \varepsilon\,\big] \leq 2\exp(-m\varepsilon^2/8)
\]
上界仅含 $m$，与 $\|x\|_\infty$ 无关。令其 $\leq \delta$ 得：
\[
m \geq \frac{8}{\varepsilon^2}\ln\frac{2}{\delta} = \Theta(\varepsilon^{-2}\log\delta^{-1})
\]
即使 $\|x\|_\infty \leq O(1/\sqrt{d})$，$m$ 仍为此量级。\hfill $\square$

\textbf{4.} \textbf{坐标采样}：直接输出 $x$ 的某几维。$x$ 平坦（$\|x\|_\infty$ 小）$\Rightarrow$ 能量均匀 $\Rightarrow$ 每维方差小 $\Rightarrow$ 少量样本即准确。

\textbf{高斯 JL}：每个输出坐标已是所有 $d$ 维的加权和 $\Rightarrow$ 天然均匀化。无论 $x$ 的 $\|x\|_\infty$ 多大或多小，每个 $(\Pi x)_i$ 的方差恒为 $1/m$。故低 $\ell_\infty$-norm 不带来额外收益。


% --- Q2 ---
\subsection*{Q2 \quad Torus 上的随机网格 LSH}

空间 $([0,n]^d, d_{\text{torus}})$，网格宽 $w$，随机平移 $s \in [0,w)^d$。
$\Delta_i = \min\{|x_i-y_i|,\, n-|x_i-y_i|\}$。
$h_s(x) = (\lfloor (x_1+s_1)/w \rfloor, \dots, \lfloor (x_d+s_d)/w \rfloor)$（torus 取模）。

\textbf{1.} 若所有 $\Delta_i < w$（不受 torus 边界干扰）。

每坐标 $i$：网格线在 $[0,w)$ 均匀分布，落入 $x_i,y_i$ 之间（长 $\Delta_i$）的概率 $= \Delta_i/w$。同桶 $\Leftrightarrow$ 网格线不落入 $\Rightarrow$ 概率 $= 1 - \Delta_i/w$。各坐标独立相乘：
\[
\boxed{\Pr[h_s(x)=h_s(y)] = \prod_{i=1}^d \left(1 - \frac{\Delta_i}{w}\right)}
\]

\textbf{2.} 由 Cauchy-Schwarz：$\sum \Delta_i \leq \sqrt{d}\sqrt{\sum \Delta_i^2} \leq \sqrt{d}\cdot r$。
用 $\prod(1-a_i) \geq 1 - \sum a_i$（$a_i\geq 0$，归纳法）：
\[
\prod(1-\tfrac{\Delta_i}{w}) \geq 1 - \tfrac{1}{w}\sum\Delta_i \geq \boxed{1 - \tfrac{\sqrt{d}\,r}{w}}
\]

\textbf{3.} 随机平移 $s$ 使 axis-aligned cubes（边长 $w$）随机放置。近点 $(r \ll w)$：各坐标差小 $\Rightarrow$ 大概率同 cube。远点：多坐标有显著差 $\Rightarrow$ cube 边界切开。

\textbf{4.} 高维局限：距离均匀分散在 $d$ 维上时，即使总距离小，乘积 $(1-\Delta_i/w)$ 中每个因子 $<1$，乘积 $\to 0$（$d$ 大时指数衰减）。

\textbf{5.} 仅一个坐标相差 $r$（$r<w$），其余全同：
\[
\boxed{\Pr = 1 - \frac{r}{w}}
\]

\textbf{6.} 每坐标差 $r/\sqrt{d}$（$r/\sqrt{d}<w$），总 torus 距离 $=r$：
\[
\boxed{\Pr = \left(1 - \frac{r}{w\sqrt{d}}\right)^d}
\]

\textbf{7.} 同为 torus 距离 $r$：(5) $P=1-r/w$ vs (6) $P \approx e^{-r\sqrt{d}/w}$。分散程度不同 $\Rightarrow$ 碰撞概率不同——不是仅取决于欧氏距离。

\textbf{8.} Axis-aligned grid LSH 受坐标方向和维数影响，不是仅依赖欧氏距离的完美 LSH。距离集中少数坐标 $\to$ 好；均匀分散 $\to$ 高维灾难。


% --- Q3 ---
\subsection*{Q3 \quad 球面上的随机超平面 LSH}

球面 $\mathbb{S}^{d-1}$，$d(x,y) = \arccos\inner{x,y}$。$h_g(x) = \operatorname{sign}\inner{g,x}$，$g\sim\calN(0,I_d)$。

\textbf{1.} $\theta = \arccos\inner{x,y}$。$(g\cdot x,\, g\cdot y)$ 是联合正态：均值 $0$，方差 $1$，协方差 $x^\top \E[gg^\top]y = \inner{x,y} = \cos\theta$。

对相关系数 $\cos\theta$ 的二元标准正态 $(X,Y)$：$\Pr[XY>0] = 1 - \theta/\pi$。
\[
\boxed{\Pr[h_g(x)=h_g(y)] = 1 - \frac{\theta}{\pi}}
\]

\textbf{2.} $g$ 定义随机超平面（法向量）。$x,y$ 同侧 $\Leftrightarrow$ 超平面不穿过它们之间的球面弧。弧长 $=\theta$，超平面方向均匀 $\Rightarrow$ 穿过概率 $=\theta/\pi$，不穿过 $=1-\theta/\pi$。

\textbf{3.} $k$ 个独立 $g_i$，$H(x) = (h_{g_1}(x),\dots,h_{g_k}(x))$：
\[
\boxed{\Pr[H(x)=H(y)] = \left(1 - \frac{\theta}{\pi}\right)^k}
\]

\textbf{4.} 近点 $(\leq\theta_1)$：$p_1 = 1-\theta_1/\pi$。远点 $(\geq\theta_2)$：$p_2 = 1-\theta_2/\pi$。
\[
\boxed{\rho = \frac{\log p_1}{\log p_2} = \frac{\log(1-\theta_1/\pi)}{\log(1-\theta_2/\pi)}}
\]
小角度近似 $\log(1-x)\approx -x$：$\rho \approx \theta_1/\theta_2$。


% --- Q4 ---
\subsection*{Q4 \quad 由范数保持推出内积保持}

\textbf{1.} 极化恒等式：$\inner{x,y} = \frac{1}{4}(\|x+y\|^2 - \|x-y\|^2)$。

同理 $\inner{\Pi x,\Pi y} = \frac{1}{4}(\|\Pi(x+y)\|^2 - \|\Pi(x-y)\|^2)$。

代入 $(1-\varepsilon)\|v\|^2 \leq \|\Pi v\|^2 \leq (1+\varepsilon)\|v\|^2$（$v\in\{x,y,x+y,x-y\}$）：
\begin{align*}
\inner{\Pi x,\Pi y} - \inner{x,y}
&\leq \frac{\varepsilon}{4}(\|x+y\|^2 + \|x-y\|^2) \\
&\geq -\frac{\varepsilon}{4}(\|x+y\|^2 + \|x-y\|^2)
\end{align*}

由平行四边形法则 $\|x+y\|^2+\|x-y\|^2 = 2(\|x\|^2+\|y\|^2)$：
\[
\boxed{|\inner{\Pi x,\Pi y} - \inner{x,y}| \leq \frac{\varepsilon}{2}(\|x\|^2+\|y\|^2) = O(\varepsilon)(\|x\|^2+\|y\|^2)}
\]

\textbf{2.} $\|x\|=\|y\|=1$ 时：\boxed{$|\inner{\Pi x,\Pi y} - \inner{x,y}| \leq \varepsilon = O(\varepsilon)$}

\textbf{3.} JL 保持 $n$ 个点的两两距离 $\Rightarrow$ 引入 $\{x_i+x_j\}$ 和 $\{x_i-x_j\}$（共 $O(n^2)$ 额外点）后，JL 对扩充集合仍保持距离 $\Rightarrow$ $\Pi$ 对 $\{x,y,x+y,x-y\}$ 满足范数保持 $\Rightarrow$ 由(1)得内积保持。维数仍为 $O(\varepsilon^{-2}\log n)$。

\textbf{4.} 单位球面上 $\|x-y\|^2 = 4\sin^2(\theta/2)$。JL 保距 $\Rightarrow$ 保 $\inner{x,y}=\cos\theta$。
由 $\arccos'(t) = -1/\sqrt{1-t^2} = -1/\sin\theta$：当 $\theta$ 不接近 $0,\pi$ 时，内积误差 $O(\varepsilon)$ 转换为角度误差 $O(\varepsilon/\sin\theta)$。


% --- Q5 ---
\subsection*{Q5 \quad 从阈值 ANN 到 Approximate Nearest Neighbor}

\textbf{1.} 最近邻距离 $d^*$。令 $k = \min\{i : \exists p, \|p-q\|\leq r_i\}$。$d^* > r_{k-1}$ 且 $d^* \leq r_k = b\,r_{k-1}$。

黑盒在半径 $r_k$ 返回 $p'$：$\|p'-q\| \leq a\cdot r_k = ab\cdot r_{k-1} < ab\cdot d^*$。
\[
\boxed{\text{近似比} \leq ab}
\]

\textbf{2.} 目标 $c$-近似 $\Rightarrow$ 需 $ab \leq c$。取 \boxed{$a = b = \sqrt{c}$}。

\textbf{3.} $a=\sqrt{c}$ 使黑盒负担更轻（$\sqrt{c} < c$），$b=\sqrt{c}$ 使半径增长率温和，层级数可控。

\textbf{4.} $r_0=1$，$r_i=b^i$。$b^K \geq \Delta \Rightarrow K \geq \log_b\Delta$。
\[
\boxed{\text{层级数} = \lceil\log_b\Delta\rceil = O\!\left(\frac{\log\Delta}{\log b}\right)}
\]


% --- Q6 ---
\subsection*{Q6 \quad 先用 JL 降维，再做随机超平面 LSH}

\textbf{1.} 已知 $|\|\Pi x-\Pi y\|^2 - \|x-y\|^2| \leq \varepsilon$ 且 $|\|\Pi x\|^2 - 1| \leq \varepsilon$。

由范数保持推内积保持（见 Q4），对 $x,y,x+y,x-y$ 范数均在 $(1\pm\varepsilon)$ 内 $\Rightarrow$
\[
\boxed{|\inner{\Pi x,\Pi y} - \inner{x,y}| \leq O(\varepsilon)}
\]

\textbf{2.} $\theta = \arccos\inner{x,y}$，$\tilde\theta = \arccos\!\big(\frac{\inner{\Pi x,\Pi y}}{\|\Pi x\|\|\Pi y\|}\big)$。

分子 $\approx \inner{x,y} \pm O(\varepsilon)$，分母 $\approx 1 \pm O(\varepsilon)$。归一化内积 $\approx \cos\theta \pm O(\varepsilon)$。

由 $\arccos'(t) = -1/\sqrt{1-t^2} = -1/\sin\theta$：
\[
\boxed{|\tilde\theta - \theta| \leq O\!\left(\frac{\varepsilon}{\sin\theta}\right)}
\]

\textbf{3.} JL 将 $d \to m = O(\varepsilon^{-2}\log n) \ll d$。在低维空间运行超平面 LSH，碰撞概率 $1-\tilde\theta/\pi \approx 1-(\theta\pm O(\varepsilon))/\pi$，时间复杂度从 $O(d)$ 降为 $O(m)$。

\textbf{4.} $\theta$ 小时 $\sin\theta \approx \theta$，误差 $O(\varepsilon/\theta)$。欲使相对误差 $\leq 50\%$（$|\tilde\theta-\theta| \leq \theta/2$）：
\[
O(\varepsilon/\theta) \leq \theta/2 \;\Rightarrow\; \boxed{\varepsilon \leq O(\theta^2)}
\]
近点需要更小的 JL 误差 $\Rightarrow$ 更大的 $m$。


% --- Q7 ---
\subsection*{Q7 \quad 为什么 CountSketch 需要随机符号}

\textbf{1.} 无符号版：$\Pi_{h(i),i}=1$。$x = (1,\dots,1)/\sqrt{n}$。$B_k = \#\{i:h(i)=k\} \sim \operatorname{Binomial}(n,1/m)$。

$\|\Pi x\|^2 = \frac{1}{n}\sum_{k=1}^m B_k^2$。$\E[B_k^2] = \Var(B_k) + \E[B_k]^2 = \frac{n}{m}(1-\frac{1}{m}) + \frac{n^2}{m^2}$。
\[
\E[\|\Pi x\|^2] = \frac{m}{n}\!\left[\frac{n}{m}\!\left(1-\frac{1}{m}\right) + \frac{n^2}{m^2}\right] = \boxed{1 + \frac{n-1}{m}}
\]

\textbf{2.} $m=o(n)$ 时 $\E[\|\Pi x\|^2] \gg 1$，真实 $\|x\|^2=1$。严重偏高。根因：碰撞 $=$ 正数相加，产生系统性正偏差。

\textbf{3.} 标准版 $\Pi_{h(i),i} = \sigma_i \in \{\pm 1\}$。桶 $k$ 的值为 $\sum_{i:h(i)=k} \sigma_i/\sqrt{n}$。

平方 $= \frac{1}{n}\big(\sum\sigma_i^2 + \sum_{i\neq j}\sigma_i\sigma_j\big) = \frac{1}{n}\big(B_k + \sum_{i\neq j}\sigma_i\sigma_j\big)$。

$\E[\sigma_i\sigma_j] = \E[\sigma_i]\E[\sigma_j] = 0$（独立，均值 $0$）。交叉项期望全部为零。
\[
\E[\|\Pi x\|^2] = \frac{1}{n}\sum_k \E[B_k] = 1 \quad\text{（无偏）}
\]

\textbf{4.} 随机符号不能避免碰撞，但使碰撞产生的交叉项 $\sigma_i\sigma_j$ 期望为零。误差从系统性偏差变为零均值随机噪声，可通过增大 $m$ 控制方差。


% --- Q8 ---
\subsection*{Q8 \quad CountSketch OSE 的 $m = \Omega(d^2)$ 下界}

$\Pi$ 每列只有一个非零元（位置和值任意确定）。随机矩阵 $A$：每列 $j\in[d]$ 独立随机选一行 $r_j\in[n]$，$A_{r_j,j}=1$。

$Ae_j =$ 第 $r_j$ 标准基向量。$\Pi A e_j = \Pi$ 的第 $r_j$ 列（仅一行非零，记为 $h(r_j)$）。

若 $h(r_i) = h(r_j)$（碰撞），$\Pi A(e_i+e_j)$ 的范数估计出错（两个非零元在同一行，可能抵消或加倍）。

为保证 OSE 成立，需对所有 $i\neq j$ 避免碰撞。

$d$ 个独立随机行号投 $m$ 个桶。无碰撞概率：
\[
\Pr[\text{无碰撞}] = \prod_{i=0}^{d-1}\!\left(1 - \frac{i}{m}\right) \approx \exp\!\left(-\frac{d^2}{2m}\right)
\]

要此概率 $\geq \Omega(1)$：$d^2/(2m) = O(1) \Rightarrow \boxed{m = \Omega(d^2)}$。

若 $m = o(d^2)$，以高概率存在碰撞 $\Rightarrow$ $\Pi$ 不是 OSE。故 column sparsity $=1$ 的 OSE 必须 $m=\Omega(d^2)$。


% --- Q9 ---
\subsection*{Q9 \quad CountSketch OSE 的 $m = O(d^2/\varepsilon^2\delta)$ 上界}

$U\in\R^{n\times d}$，$U^\top U = I_d$。$\Pi$：$h(i)\in[m]$，$\sigma_i\in\{\pm1\}$，$\Pi_{h(i),i}=\sigma_i$。
$u_i = U$ 的第 $i$ 行（$d$ 维列向量）。

\textbf{1.} $(\Pi^\top\Pi)_{i,j} = \sum_{k=1}^m \Pi_{k,i}\Pi_{k,j}$。
$i=j$：$=\sigma_i^2=1 \Rightarrow (\Pi^\top\Pi - I_n)_{i,i}=0$。
$i\neq j$：仅当 $h(i)=h(j)$ 时为 $\sigma_i\sigma_j$。
\[
\boxed{(\Pi^\top\Pi - I_n)_{i,j} = \ind{h(i)=h(j)}\sigma_i\sigma_j \;\;(i\neq j),\quad\text{对角元}=0}
\]

\textbf{2.}
$U^\top\Pi^\top\Pi U - I_d = U^\top(\Pi^\top\Pi - I_n)U$（因 $U^\top U = I_d$）
\[
\boxed{= \sum_{i\neq j:\,h(i)=h(j)} \sigma_i\sigma_j\, u_i u_j^\top}
\]

\textbf{3.} Frobenius 平方展开为双重和，含因子 $\sigma_i\sigma_j\sigma_k\sigma_\ell$。
$\E[\sigma_i]=0$ 独立 $\Rightarrow$ 若 $\{i,j\}\neq\{k,\ell\}$ 则某 $\sigma$ ``独苗''，期望为零。
仅 $\{i,j\}=\{k,\ell\}$（自身对）时 $\E[\sigma_i^2\sigma_j^2]=1$ 保留。

\textbf{4.} 保留项：$\sum_{i\neq j} \Pr[h(i)=h(j)]\cdot\|u_i u_j^\top\|_F^2$。
$\Pr[h(i)=h(j)] = 1/m$。$\|u_i u_j^\top\|_F^2 = \|u_i\|^2\|u_j\|^2$。
\[
\boxed{\E[\|U^\top\Pi^\top\Pi U - I_d\|_F^2] \leq \frac{1}{m}\sum_{i\neq j}\|u_i\|^2\|u_j\|^2}
\]

\textbf{5.} $\sum_i \|u_i\|^2 = \|U\|_F^2 = d$。$\sum_{i\neq j}\|u_i\|^2\|u_j\|^2 \leq (\sum_i\|u_i\|^2)^2 = d^2$。
\[
\boxed{\E[\|U^\top\Pi^\top\Pi U - I_d\|_F^2] \leq \frac{d^2}{m}}
\]

\textbf{6.} Markov：$\Pr[\|M\|_F > \varepsilon] = \Pr[\|M\|_F^2 > \varepsilon^2] \leq \E[\|M\|_F^2]/\varepsilon^2 \leq d^2/(m\varepsilon^2) \leq \delta$。
\[
\boxed{m \geq \frac{d^2}{\varepsilon^2\delta} \;\Rightarrow\; \|M\|_F \leq \varepsilon \text{ w.p. } \geq 1-\delta}
\]

\textbf{7.} $\forall x\in\R^d, y=Ux$：
$|\|\Pi Ux\|^2 - \|Ux\|^2| = |x^\top M x| \leq \|M\|_{\text{op}}\|x\|^2 \leq \|M\|_F\|x\|^2 \leq \varepsilon\|x\|^2$。
又 $\|Ux\|^2 = \|x\|^2$，故
\[
\boxed{(1-\varepsilon)\|Ux\|^2 \leq \|\Pi Ux\|^2 \leq (1+\varepsilon)\|Ux\|^2}
\]

\textbf{8.} CountSketch 每列仅一个非零元，取 $m=O(d^2/(\varepsilon^2\delta))$ 即为 $d$ 维子空间的 $\varepsilon$-OSE。结合 Q8 下界 $\Omega(d^2)$，\textbf{$m=\Theta(d^2)$ 是 column sparsity $=1$ OSE 的紧界}。


% ============================================================
\section*{Ch5：网络流与线性规划（Flow \& LP）}
% ============================================================

% --- Q10 ---
\subsection*{Q10 \quad 点容量最大流}

\textbf{建图（拆点法）}：

对每个 $v \in V\setminus\{s,t\}$：拆为 $v_{\text{in}}, v_{\text{out}}$，加边 $(v_{\text{in}} \to v_{\text{out}})$，容量 $= c_v$。

对每条原边 $(u,v)$：加边 $(u_{\text{out}} \to v_{\text{in}})$，容量 $= \infty$（$s,t$ 直接连）。

\textbf{正确性}：任何流经 $v$ 的流量必过 $(v_{\text{in}}\to v_{\text{out}})$ 且 $\leq c_v$，对应点容量约束；新图中流与原图可行流一一对应且保流值。

\textbf{算法}：新图上 Dinic/Edmonds-Karp，$\leq 2n$ 点，$\leq m+n$ 边，多项式时间。


% --- Q11 ---
\subsection*{Q11 \quad 二分图最大匹配与最小点覆盖的 LP}

\textbf{最大匹配 ILP}：
\[
\max \sum_{e\in E} x_e \quad\text{s.t.}\quad \sum_{e\ni v} x_e \leq 1\;( \forall v),\; x_e \in \{0,1\}
\]
\textbf{LP 松弛}：$x_e \geq 0$。

\textbf{对偶 LP}：
\[
\min \sum_{v\in V} y_v \quad\text{s.t.}\quad y_u + y_v \geq 1\;(\forall (u,v)\in E),\; y_v \geq 0
\]

\textbf{最小点覆盖 ILP}：$\min\sum y_v$，s.t. $y_u+y_v\geq 1$，$y_v\in\{0,1\}$。其 LP 松弛 $=$ 上述对偶。

\textbf{证明 $\nu(G) = \tau(G)$}：

(1) 二分图关联矩阵全单模 $\Rightarrow$ 匹配 LP 基本可行解为整数 $\Rightarrow$ LP$^* = \nu(G)$。

(2) 对偶 LP 约束矩阵亦全单模 $\Rightarrow$ 对偶 LP$^* = \tau(G)$。

(3) 强对偶性：原始 LP$^* =$ 对偶 LP$^*$。故 \boxed{$\nu(G) = \tau(G)$}（K\H{o}nig 定理）。


% --- Q12 ---
\subsection*{Q12 \quad 背包问题的 LP 松弛与取整}

\textbf{ILP}：$\max\sum v_i x_i$，s.t. $\sum w_i x_i \leq C$，$x_i\in\{0,1\}$。
\textbf{LP}：放松为 $0\leq x_i\leq 1$，记最优解 $x^*$。

\textbf{取整算法}：
(1) 物品按 $v_i/w_i$ 降序排列。
(2) 贪心取直到装不下。记贪心值 $V_{\text{greedy}}$，设 $k$ 为第一装不下物品。
(3) 输出 $\max(V_{\text{greedy}}, v_k)$。

\textbf{近似比证明}：

OPT$_{\text{LP}} = \sum v_i x^*_i \geq$ OPT。

$V_{\text{greedy}} + v_k \geq \sum_{i=1}^k v_i x^*_i \geq$ OPT$_{\text{LP}} \geq$ OPT。

故 $\max(V_{\text{greedy}}, v_k) \geq$ OPT$/2$。\boxed{\text{1/2-近似}}。


% --- Q13 ---
\subsection*{Q13 \quad 带下界的可行流}

\textbf{算法}：

(1) 每边新容量 $c'_e = u_e - \ell_e$，强制预流 $\ell_e$。

(2) 每点净需求：$d(v) = \sum_{\text{入}} \ell_e - \sum_{\text{出}} \ell_e$。

(3) 加超级源 $S^*$、汇 $T^*$：
\begin{itemize}
\item $d(v) > 0$（净流入）$\to$ 加边 $(S^* \to v)$，容量 $= d(v)$
\item $d(v) < 0$（净流出）$\to$ 加边 $(v \to T^*)$，容量 $= -d(v)$
\end{itemize}

(4) 求 $S^*\to T^*$ 最大流 $f'$。

\textbf{判据}：\boxed{\text{存在可行流} \;\Leftrightarrow\; S^*\text{ 的所有出边被 } f' \text{ 饱和}}。

\textbf{构造}：可行流 $f_e = \ell_e + f'_e$。

\textbf{正确性}：预流 $\ell_e$ 造成不平衡 $d(v)$，$f'$ 校正此不平衡使每点流量守恒。$S^*$ 出边饱和 $\Leftrightarrow$ 所有净需求被满足。


% --- Q14 ---
\subsection*{Q14 \quad DAG 最小路径覆盖}

\textbf{归约二分图}：构造 $H=(L,R,E')$：
\begin{itemize}
\item $L=\{\ell_1,\dots,\ell_n\}$（出发副本），$R=\{r_1,\dots,r_n\}$（到达副本）
\item $\forall (u,v)\in E$：加边 $(\ell_u, r_v) \in E'$
\end{itemize}

求 $H$ 的最大匹配 $M$。\boxed{\text{最小路径覆盖} = |V| - M}。

\textbf{证明}：匹配边 $(\ell_u, r_v) \Leftrightarrow$ 路径中 $u$ 紧接 $v$。
初始 $n$ 条单点路径，每匹配边拼接两条路径片段 $\Rightarrow$ 路径数减 $1$。
最大匹配 $\Rightarrow$ 最小路径数。


% --- Q15 ---
\subsection*{Q15 \quad 点不相交路径}

\textbf{拆点法}：

对 $v \in V\setminus\{s,t\}$：拆 $v_{\text{in}},v_{\text{out}}$，加边 $(v_{\text{in}}\to v_{\text{out}})$，容量 $= \bm{1}$。

原边 $(u,v)$：$(u_{\text{out}}\to v_{\text{in}})$，容量 $=1$（$s,t$ 直接连）。

\textbf{求 $s\to t$ 最大流}，流值 $=$ 最大内部点不相交路径数。

\textbf{Menger 定理（点版本）}：最大内部点不相交 $s$-$t$ 路径数 $=$ 最小 $s$-$t$ 点割大小。


% --- Q16 ---
\subsection*{Q16 \quad 带权点覆盖的 Primal-Dual 算法}

\textbf{ILP}：$\min\sum w_v x_v$，s.t. $x_u+x_v\geq 1$，$x_v\in\{0,1\}$。
\textbf{LP}：$x_v \geq 0$。

\textbf{对偶 LP}：$\max\sum_{e\in E} y_e$，s.t. $\sum_{e\ni v} y_e \leq w_v$，$y_e \geq 0$。

\textbf{Primal-Dual 算法}：

(1) 初始化 $y_e=0$，$x_v=0$，所有边未覆盖。

(2) WHILE 存在未覆盖边 $(u,v)$（$x_u=x_v=0$）：
持续增大 $y_{uv}$ 直到 $\exists$ 端点满足 $\sum_{e\ni v} y_e = w_v$。
将满足条件的端点选入覆盖（$x_v=1$），邻边标已覆盖。

(3) 返回 $C = \{v: x_v=1\}$。

\textbf{近似比 $=2$ 的证明}：

(1) 算法始终维护 $\sum_{e\ni v}y_e \leq w_v$ $\Rightarrow$ $y$ 是对偶可行解。

(2) $\sum_{v\in C} w_v = \sum_{v\in C}\sum_{e\ni v} y_e$（选中时约束紧）。
每条边 $e$ 的 $y_e$ 最多被两端点各计一次 $\leq 2y_e$。
故 $\sum_{v\in C} w_v \leq 2\sum_{e\in E} y_e$。

(3) 弱对偶 $\sum y_e \leq$ OPT$_ {\text{LP}} \leq$ OPT $\Rightarrow$ \boxed{$\sum_{v\in C} w_v \leq 2\cdot\text{OPT}$}。


% --- Q17 ---
\subsection*{Q17 \quad Set Cover 的 LP 松弛与随机取整}

\textbf{ILP}：$\min\sum c_S x_S$，s.t. $\sum_{S\ni e} x_S \geq 1$，$x_S\in\{0,1\}$。
\textbf{LP}：$x_S \geq 0$，记最优解 $x^*_S$。

\textbf{随机取整算法}：每个 $S$ 独立以概率 $x'_S = \min\{1,\, \alpha x^*_S\}$ 选择，$\alpha = \Theta(\log|U|)$。

\textbf{覆盖概率}：对元素 $e$，
\begin{align*}
\Pr[e\text{ 未被覆盖}] &= \prod_{S\ni e}(1-x'_S)
\leq \prod_{S\ni e} e^{-\alpha x^*_S}
= \exp\!\big(-\alpha\!\sum_{S\ni e} x^*_S\big) \\
&\leq e^{-\alpha} \qquad (\text{LP 约束 } \sum x^*_S \geq 1)
\end{align*}

取 $\alpha = c\ln|U|$（$c>1$）：$\Pr[e\text{ 未覆盖}] \leq |U|^{-c}$。

Union bound：$\Pr[\exists\text{未覆盖}] \leq |U|^{1-c} \to 0$（$c\geq2$）。

\textbf{期望成本}：$\E[\text{成本}] = \sum c_S x'_S \leq \alpha\sum c_S x^*_S = \alpha\cdot\text{OPT}_{\text{LP}} \leq \alpha\cdot\text{OPT}$。

取 $\alpha = \Theta(\log|U|)$：\boxed{O(\log|U|)\text{-近似}}。


% --- Q18 ---
\subsection*{Q18 \quad 带点容量和需求的最小费用运输问题}

\textbf{建图}：

节点：超级源 $S$、超级汇 $T$；工厂 $F_i$；仓库拆 $W_v^{\text{in}}, W_v^{\text{out}}$；商店 $R_j$。

边：
\begin{itemize}
\item $(S \to F_i)$：容量 $= s_i$（供应上限），费用 $= 0$
\item $(F_i \to W_v^{\text{in}})$：容量 $= u_e$，费用 $= a_e$
\item $(W_v^{\text{in}} \to W_v^{\text{out}})$：容量 $= c_v$（仓库吞吐），费用 $= 0$
\item $(W_v^{\text{out}} \to R_j)$：容量 $= u_e$，费用 $= a_e$
\item $(R_j \to T)$：容量 $= d_j$（需求），费用 $= 0$
\end{itemize}

\textbf{求解}：$S\to T$ 最小费用最大流（Successive Shortest Path）。

\textbf{判据}：最大流 $= \sum_j d_j$ $\Leftrightarrow$ 可行；最小费用即最优总成本。

\textbf{正确性}：$S\to F_i$ 限供 $\leq s_i$；拆点 $c_v$ 限仓库流量 $\leq c_v$；$R_j\to T$ 容量 $d_j$ 且饱和 $\Rightarrow$ 需求满足；流量守恒 $\Rightarrow$ 物理可行性。目标函数 $\sum a_e f_e$ 即总成本。


% ============================================================
\vspace{1em}
\begin{center}
\rule{0.6\textwidth}{0.4pt}\\[4pt]
{\large \textbf{18 题考场速写版完成 · 按教师通知的 2 题背诵}}
\end{center}

\end{document}
